\section*{Capítulo \ref{ch:b1:diffeq} --- Ecuaciones diferenciales lineales}

\begin{solution}{exo:b1:diffeq:1}
Homogénea: $y_h = \lambda\,\eu^{-2x}$. Particular: pruébese
$y_p = c\,\eu^{3x}$ ($3$ no es raíz de $r + 2$): $3c + 2c = 1$,
$c = \frac15$. Solución general: $y = \frac{\eu^{3x}}{5} +
\lambda\,\eu^{-2x}$. Con $y(0) = 1$: $\frac15 + \lambda = 1$,
$\lambda = \frac45$, luego $y = \frac{\eu^{3x} + 4\,\eu^{-2x}}{5}$.
\end{solution}

\begin{solution}{exo:b1:diffeq:2}
En $\intoo{0}{+\infty}$, divídase por $x$: $y' - \frac{2}{x}\,y = x^2$.
Aquí $A(x) = -2\ln x$, $\eu^{-A(x)} = x^2$: soluciones homogéneas
$\lambda x^2$. \omterm{thm:b1:diffeq:voc}{Variación de las constantes}: $\mu'(x) = x^2 \cdot x^{-2} =
1$, luego $\mu = x + \lambda$ y
\[
y(x) = x^3 + \lambda x^2, \qquad \lambda \in \R .
\]
\emph{Comprobación:} $x(3x^2 + 2\lambda x) - 2(x^3 + \lambda x^2) = x^3$.
\end{solution}

\begin{solution}{exo:b1:diffeq:3}
$y'' - 3y' + 2y = 0$: raíces $1$ y $2$; $\;y = \lambda\,\eu^{x} +
\mu\,\eu^{2x}$.

$y'' + 4y' + 4y = 0$: raíz doble $-2$; $\;y = (\lambda + \mu
x)\,\eu^{-2x}$.

$y'' - 2y' + 5y = 0$: raíces $1 \pm 2\iu$; $\;y = \eu^{x}(\lambda\cos
2x + \mu\sin 2x)$.
\end{solution}

\begin{solution}{exo:b1:diffeq:4}
Homogénea: raíces $\pm 1$, $y_h = \lambda\,\eu^x + \mu\,\eu^{-x}$.
Particular con miembro derecho polinómico ($\gamma = 0$ no es raíz):
$y_p = ax^2 + bx + c$; sustituyendo, $2a - (ax^2 + bx + c) = x^2$ da
$a = -1$, $b = 0$, $c = 2a = -2$: $y_p = -x^2 - 2$. Solución general
$y = -x^2 - 2 + \lambda\eu^x + \mu\eu^{-x}$.

Cauchy: $y(0) = -2 + \lambda + \mu = 0$ e $y'(0) = \lambda - \mu = 1$:
$\lambda = \frac32$, $\mu = \frac12$. Luego
$y = -x^2 - 2 + \frac{3\eu^x + \eu^{-x}}{2}$.
\end{solution}

\begin{solution}{exo:b1:diffeq:5}
$a(x) = \tan x$, $A(x) = -\ln(\cos x)$ (válido: $\cos > 0$ en el
intervalo), $\eu^{-A} = \cos x$: soluciones homogéneas $\lambda\cos x$.
\omterm{thm:b1:diffeq:voc}{Variación de las constantes}: $\mu'(x) = \sin 2x \cdot \frac{1}{\cos x} =
2\sin x$, luego $\mu = -2\cos x + \lambda$ y
\[
y(x) = -2\cos^2 x + \lambda \cos x .
\]
\emph{Comprobación:} $y' = 4\cos x \sin x - \lambda\sin x$ y
$y\tan x = -2\cos x\sin x + \lambda \sin x$; su suma es
$2\cos x\sin x = \sin 2x$, como se pedía.
\end{solution}

\begin{solution}{exo:b1:diffeq:6}
$\chi(r) = r^2 - 4r + 3 = (r-1)(r-3)$: $\gamma = 1$ es raíz simple
($m = 1$). Pruébese $y_p = x(ax + b)\,\eu^x$. Con $u = ax^2 + bx$,
\[
y_p'' - 4y_p' + 3y_p = \bigl(u'' + (2 - 4)u' + \chi(1) u\bigr)\eu^x
= \bigl(2a - 2(2ax + b)\bigr)\eu^x .
\]
Identifíquese con $(2x + 1)\eu^x$: $-4a = 2$ y $2a - 2b = 1$, luego
$a = -\frac12$, $b = -1$. Solución general:
\[
y = -\Bigl(\frac{x^2}{2} + x\Bigr)\eu^{x} + \lambda\,\eu^{x} +
\mu\,\eu^{3x}, \qquad (\lambda, \mu) \in \R^2 .
\]
\end{solution}

\begin{solution}{exo:b1:diffeq:7}
Homogénea: $y_h = \lambda\cos 2x + \mu\sin 2x$.

Miembro derecho $x$ ($\gamma = 0$ no es raíz): $y_1 = ax + b$ con
$4(ax + b) = x$: $y_1 = \frac x4$.

Miembro derecho $\sin 2x = \Im(\eu^{2\iu x})$, con $2\iu$ raíz simple
de $r^2 + 4$: pruébese $z = c\,x\,\eu^{2\iu x}$; entonces
$z'' + 4z = 4\iu c\,\eu^{2\iu x}$, igual a $\eu^{2\iu x}$ para
$c = \frac{1}{4\iu} = -\frac{\iu}{4}$. Luego $z = -\frac{\iu x}{4}(\cos 2x +
\iu \sin 2x)$ e $y_2 = \Im(z) = -\frac{x\cos 2x}{4}$.

Por superposición:
\[
y = \frac{x}{4} - \frac{x\cos 2x}{4} + \lambda\cos 2x + \mu\sin 2x .
\]
\end{solution}

\begin{solution}{exo:b1:diffeq:8}
La ecuación $T' + kT = 20k$ tiene la solución particular constante $20$
y soluciones homogéneas $\lambda\eu^{-kt}$: $T(t) = 20 +
\lambda\,\eu^{-kt}$, y $T(0) = 80$ da $\lambda = 60$:
\[
T(t) = 20 + 60\,\eu^{-kt} .
\]
$T(10) = 50$: $\eu^{-10k} = \frac12$, luego $k = \frac{\ln 2}{10}$.
Después, $T(t) = 25$ exige $\eu^{-kt} = \frac{5}{60} = \frac{1}{12}$, es decir,
\[
t = \frac{\ln 12}{k} = 10\,\frac{\ln 12}{\ln 2} \approx 35.8
\text{ minutos.}
\]
\end{solution}

\begin{solution}{exo:b1:diffeq:9}
\begin{enumerate}
  \item $\chi(r) = r^2 + 2\varepsilon r + 1$, $\Delta = 4(\varepsilon^2
        - 1)$.
        Para $\varepsilon \in \intco{0}{1}$: raíces $-\varepsilon \pm
        \iu\sqrt{1 - \varepsilon^2}$, luego
        $y = \eu^{-\varepsilon t}\bigl(\lambda\cos\omega t +
        \mu\sin\omega t\bigr)$ con $\omega = \sqrt{1 -
        \varepsilon^2}$.
        Para $\varepsilon = 1$: raíz doble $-1$, $y = (\lambda + \mu
        t)\,\eu^{-t}$.
        Para $\varepsilon > 1$: raíces reales $r_\pm = -\varepsilon \pm
        \sqrt{\varepsilon^2 - 1}$, ambas $< 0$, y $y =
        \lambda\eu^{r_+t} + \mu\eu^{r_-t}$.
  \item Para $\varepsilon \in \intoo{0}{1}$: $\abs y \leq
        \eu^{-\varepsilon t}(\abs\lambda + \abs\mu) \to 0$. Para
        $\varepsilon = 1$: $(\lambda + \mu t)\eu^{-t} \to 0$ (la
        exponencial gana al polinomio,
        \cref{prop:b1:functions:powerrules}). Para $\varepsilon > 1$:
        las dos exponenciales decaen, pues $r_\pm < 0$ (en efecto,
        $\sqrt{\varepsilon^2 - 1} < \varepsilon$). Para
        $\varepsilon = 0$: $y = \lambda\cos t + \mu\sin t$ tiene
        amplitud constante $\sqrt{\lambda^2 + \mu^2} \neq 0$ salvo si
        $y = 0$.
  \item Escríbase $\lambda\cos\omega t + \mu\sin\omega t =
        R\cos(\omega t - \varphi)$ con $R = \sqrt{\lambda^2 + \mu^2}
        > 0$. Los ceros de $y$ son los de $\cos(\omega t -
        \varphi)$ (el factor $\eu^{-\varepsilon t}$ no se anula
        nunca): $\omega t - \varphi \equiv \frac\pi2 \pmod \pi$, una
        progresión aritmética de paso $\frac{\pi}{\omega} =
        \frac{\pi}{\sqrt{1 - \varepsilon^2}}$.
\end{enumerate}
\end{solution}

\begin{solution}{exo:b1:diffeq:10}
Hágase $x = y = 0$: $2f(0) = 2f(0)^2$, y $f(0) \neq 0$ da $f(0) = 1$.
Fíjese $x$ y derívese dos veces la ecuación respecto de $y$:
\[
f''(x+y) + f''(x-y) = 2 f(x) f''(y) .
\]
Haciendo $y = 0$: $\;2f''(x) = 2 f(x) f''(0)$, es decir,
\[
f''(x) = c\,f(x), \qquad c = f''(0).
\]
\emph{Caso $c = \omega^2 > 0$:} $f(x) = \lambda\cosh\omega x +
\mu\sinh\omega x$; $f(0) = 1$ da $\lambda = 1$. Al llevarlo a la
ecuación funcional y usar las fórmulas de adición
(\cref{prop:b1:functions:hyprules}), la ecuación obliga a $\mu = 0$
(compárense los coeficientes de $\sinh\omega x \sinh\omega y$, o
evalúese en $x = y$): $f = \cosh\omega x$, que en efecto cumple
$\cosh(x+y) + \cosh(x-y) = 2\cosh x\cosh y$.

\emph{Caso $c = -\omega^2 < 0$:} análogamente, $f(x) = \cos\omega x$
($\omega \neq 0$), que cumple la ecuación.

\emph{Caso $c = 0$:} $f$ afín con $f(0) = 1$: $f(x) = 1 + \mu x$; la
ecuación obliga a $\mu = 0$, excluido ($f$ no es constante).

Conclusión: las soluciones son $f(x) = \cos\omega x$ y
$f(x) = \cosh\omega x$, con $\omega > 0$.
\end{solution}

\begin{solution}{exo:b1:diffeq:11}
Póngase $z(t) = y(\eu^t)$, de modo que $y(x) = z(\ln x)$ para $x > 0$. Entonces
\[
y'(x) = \frac{z'(\ln x)}x,
\qquad
y''(x) = \frac{z''(\ln x) - z'(\ln x)}{x^2} ,
\]
y sustituyendo en la ecuación:
\[
x^2y'' - xy' + y = \bigl(z'' - z'\bigr) - z' + z
= z'' - 2z' + z = 0 .
\]
\omterm{def:b1:diffeq:linear2}{Polinomio característico} $(r - 1)^2$: raíz doble $1$, luego
$z(t) = (\lambda + \mu t)\,\eu^t$ y, de vuelta a la variable $x = \eu^t$:
\[
y(x) = (\lambda + \mu\ln x)\,x,
\qquad \lambda, \mu \in \R .
\]
\end{solution}

\begin{solution}{exo:b1:diffeq:12}
\begin{enumerate}
  \item En forma normalizada $y' - \frac2x\,y = 0$ en cada intervalo:
        $A(x) = -2\ln\abs x$, luego las soluciones son $y = a x^2$ en
        $\intoo0{+\infty}$ e $y = b x^2$ en $\intoo{-\infty}0$, con
        constantes independientes
        (\cref{thm:b1:diffeq:homogeneous1}).
  \item Sea $y = ax^2$ para $x \geq 0$ y $bx^2$ para $x < 0$. En cada
        semirrecta abierta, $y$ es derivable y cumple $xy' = 2y$. En
        $0$: los cocientes incrementales $\frac{y(h) - y(0)}h = ah$ o
        $bh$ tienden a $0$, luego $y'(0) = 0$ existe, y la ecuación en
        $x = 0$ dice $0 \cdot y'(0) = 2y(0) = 0$: se cumple. Así pues,
        $y$ resuelve la ecuación en todo $\R$.
  \item Las soluciones en $\R$ son exactamente estas funciones
        pegadas: una familia de \emph{dos} parámetros para una
        ecuación de primer orden. No hay contradicción con
        el~\cref{thm:b1:diffeq:voc}, cuyas hipótesis fallan aquí:
        escrita como $y' + a(x)y = 0$, el coeficiente
        $a(x) = -\frac2x$ no es continuo en $0$ --- ni siquiera está
        definido ---, de modo que $\R$ no es un intervalo donde el
        teorema se aplique. La singularidad en $0$ desconecta las dos
        semirrectas, y el valor $y(0) = 0$ queda forzado sin
        transportar información de un lado a otro. Todo dato de Cauchy
        en $x_0 \neq 0$ determina la solución solo en la semirrecta
        que contiene a $x_0$.
\end{enumerate}
\end{solution}

\begin{solution}{pb:b1:diffeq:1}
\textbf{1.} $\chi(r) = r^2 + 2\lambda r + \omega_0^2$,
$\Delta = 4(\lambda^2 - \omega_0^2) < 0$ para $0 < \lambda < \omega_0$:
raíces $-\lambda \pm \iu\omega_d$, luego, por
el~\cref{thm:b1:diffeq:homogeneous2},
\[
x(t) = \eu^{-\lambda t}\bigl(\lambda_1\cos\omega_d t +
\mu_1\sin\omega_d t\bigr), \qquad (\lambda_1, \mu_1) \in \R^2 .
\]
Para $\lambda = 0$: $x = \lambda_1\cos\omega_0 t + \mu_1\sin\omega_0 t$.
Los regímenes crítico ($\lambda = \omega_0$) y sobreamortiguado
($\lambda > \omega_0$) son los del~\cref{exo:b1:diffeq:9} (tras
reescalar el tiempo): $(\lambda_1 + \mu_1 t)\eu^{-\lambda t}$, resp.\
combinaciones de $\eu^{r_\pm t}$ con
$r_\pm = -\lambda \pm \sqrt{\lambda^2 - \omega_0^2}$.

\textbf{2.} Subamortiguado: $\abs x \leq \eu^{-\lambda
t}(\abs{\lambda_1} + \abs{\mu_1}) \to 0$. Crítico:
$(\lambda_1 + \mu_1 t)\eu^{-\lambda t} \to 0$, pues las exponenciales
ganan a los polinomios (\cref{prop:b1:functions:powerrules}).
Sobreamortiguado: $r_- < r_+ = -\lambda + \sqrt{\lambda^2 -
\omega_0^2} < 0$ porque $\sqrt{\lambda^2 - \omega_0^2} < \lambda$; las
dos exponenciales decaen.

\textbf{3.} A lo largo de una solución de $(H)$, usando
$x'' = -2\lambda x' - \omega_0^2 x$:
\[
\mathcal E'(t) = x'x'' + \omega_0^2 x x'
= x'\bigl(-2\lambda x' - \omega_0^2 x\bigr) + \omega_0^2 xx'
= -2\lambda\,x'^2 \leq 0 .
\]
Si $x(t_0) = x'(t_0) = 0$, entonces $\mathcal E(t_0) = 0$; $\mathcal E$
es no negativa y no creciente, luego $\mathcal E \equiv 0$ en
$\intco{t_0}{+\infty}$, lo que obliga a $x \equiv 0$ ahí; para
$t \leq t_0$ se repite el argumento con $\tilde x(t) = x(2t_0 - t)$,
que resuelve la ecuación con amortiguamiento $-\lambda$ pero cumple
$\tilde{\mathcal E}(t_0) = 0$ y $\tilde{\mathcal E}' =
+2\lambda\tilde x'^2 \geq 0$ con $\tilde{\mathcal E} \geq 0$; ser no
negativa, no decreciente y nula en el extremo derecho de
$\intoc{-\infty}{t_0}$ significa ser nula en todo el intervalo. Luego
$x \equiv 0$ en $\R$ --- una demostración energética de la unicidad,
válida para todo $\lambda \geq 0$.

\textbf{4.} $x(t + T_d) = R\,\eu^{-\lambda t}\eu^{-\lambda
T_d}\cos(\omega_d t + 2\pi - \varphi) = \eu^{-\lambda T_d}x(t)$. El
factor de encogimiento por seudoperíodo es $\eu^{-\delta}$ con
$\delta = \lambda T_d = \frac{2\pi\lambda}{\omega_d}$. Para
$\omega_0 = 1$, $\lambda = 0.1$: $\omega_d = \sqrt{0.99} = 0.99499$,
luego $\delta = \frac{0.62832}{0.99499} = 0.6315$: cada oscilación
conserva $\eu^{-0.63} \approx 53\,\%$ de su amplitud.

\textbf{5.} El factor de amplitud es $\eu^{-\lambda t}$, que se divide
por $\eu$ en el tiempo $t = \frac1\lambda$. Ese intervalo contiene
$\frac{1/\lambda}{T_d} = \frac{\omega_d}{2\pi\lambda}$ seudoperíodos.
Para $\lambda \ll \omega_0$, $\omega_d \approx \omega_0$ y eso vale
$\approx \frac{\omega_0}{2\pi\lambda} = \frac Q\pi$. Una cuerda de
guitarra con $Q = 300$ suena unos cien períodos; un amortiguador de
puerta con $Q = 1$ no completa ni uno.

\textbf{6.} Sustituyendo $z\,\eu^{\iu\Omega t}$ en el miembro izquierdo
se obtiene $z\,(-\Omega^2 + 2\iu\lambda\Omega +
\omega_0^2)\,\eu^{\iu\Omega t}$, que es igual a
$A\,\eu^{\iu\Omega t}$ exactamente para $z = \frac{A}{\omega_0^2 -
\Omega^2 + 2\iu\lambda\Omega}$ (el denominador no es nulo: su parte
imaginaria es $2\lambda\Omega > 0$). Como los coeficientes son reales,
la parte real $x_p = \Re\bigl(z\eu^{\iu\Omega t}\bigr)$ resuelve la
ecuación con miembro derecho $\Re\bigl(A\eu^{\iu\Omega t}\bigr) =
A\cos\Omega t$.

\textbf{7.} Escríbase $\omega_0^2 - \Omega^2 + 2\iu\lambda\Omega =
\sqrt{D}\,\eu^{\iu\varphi}$ con $D = (\omega_0^2 - \Omega^2)^2 +
4\lambda^2\Omega^2$ y $\varphi \in \intoo0\pi$ (la parte imaginaria
$2\lambda\Omega$ es positiva), de modo que $\tan\varphi =
\frac{2\lambda\Omega}{\omega_0^2 - \Omega^2}$. Entonces
$z = \frac{A}{\sqrt D}\eu^{-\iu\varphi}$ y
\[
x_p(t) = \Re\Bigl(\frac A{\sqrt D}\,\eu^{\iu(\Omega t -
\varphi)}\Bigr) = R\cos(\Omega t - \varphi),
\qquad R = \frac A{\sqrt D} .
\]

\textbf{8.} Cuando $\Omega \to 0^+$: $D \to \omega_0^4$, luego
$R \to A/\omega_0^2$ y $\tan\varphi \to 0^+$ con
$\varphi \in \intoo0{\frac\pi2}$: $\varphi \to 0$. La masa sigue a la
fuerza de forma cuasiestática, desplazada en fuerza partido por
rigidez. Cuando $\Omega \to +\infty$: $D \sim \Omega^4$, luego
$R \sim A/\Omega^2 \to 0$, y $\varphi \to \pi$ (el complejo
$\omega_0^2 - \Omega^2 + 2\iu\lambda\Omega$ se va al segundo cuadrante
con argumento $\to \pi$): la masa apenas se mueve, y en oposición de
fase --- domina la inercia.

\textbf{9.} Por el~\cref{thm:b1:diffeq:structure2}~(1), toda solución
es $x = x_p + x_h$ con $x_h$ solución de $(H)$; por la pregunta 2,
$x_h(t) \to 0$, luego $x(t) - x_p(t) \to 0$: todas las soluciones
convergen al mismo régimen permanente. Las condiciones iniciales solo
dan forma al transitorio.

\textbf{10.} Si $x$ es una solución periódica, $x - x_p = x_h$ es una
solución periódica de $(H)$ que tiende a $0$ en $+\infty$; y una
función periódica con límite $0$ es idénticamente $0$ (sus valores en
un período se repiten para siempre, de modo que todo valor es límite de
una subsucesión que tiende a $0$). Luego $x = x_p$.

\textbf{11.} Aquí $\lambda = 1$, $\omega_0^2 = 2$, $\Omega = 1$,
$A = 1$: $z = \frac1{2 - 1 + 2\iu} = \frac{1 - 2\iu}5$, luego
\[
x_p = \Re\Bigl(\frac{(1 - 2\iu)(\cos t + \iu\sin t)}5\Bigr)
= \frac{\cos t + 2\sin t}5 .
\]
Homogénea: las raíces de $r^2 + 2r + 2$ son $-1 \pm \iu$:
$x_h = \eu^{-t}(C\cos t + S\sin t)$. Condiciones:
$x(0) = \frac15 + C = 0$ da $C = -\frac15$; derivando,
$x'(0) = \frac25 - C + S = 0$ da $S = C - \frac25 = -\frac35$. Por tanto
\[
x(t) = \underbrace{\frac{\cos t + 2\sin t}5}_{\text{permanente}}
- \underbrace{\eu^{-t}\,\frac{\cos t + 3\sin
t}5}_{\text{transitorio}} ,
\]
con el transitorio muriendo como $\eu^{-t}$.

\textbf{12.} Desarrollando, $g(u) = u^2 - 2(\omega_0^2 - 2\lambda^2)u
+ \omega_0^4 = (u - u_r)^2 + g(u_r)$ con $u_r = \omega_0^2 -
2\lambda^2$ y
\[
g(u_r) = \omega_0^4 - u_r^2 = (\omega_0^2 - u_r)(\omega_0^2 +
u_r) = 2\lambda^2\bigl(2\omega_0^2 - 2\lambda^2\bigr) =
4\lambda^2(\omega_0^2 - \lambda^2) .
\]
Si $2\lambda^2 < \omega_0^2$, entonces $u_r > 0$ es un cuadrado de
frecuencia admisible: $g$ tiene ahí un mínimo estricto, luego
$R = A/\sqrt g$ tiene un máximo estricto en
$\Omega_r = \sqrt{u_r} = \sqrt{\omega_0^2 - 2\lambda^2}$, con
$R_{\max} = A/\sqrt{g(u_r)} = \frac{A}{2\lambda\sqrt{\omega_0^2 -
\lambda^2}}$.

\textbf{13.} $R(0) = A/\omega_0^2$, luego
\[
\frac{R_{\max}}{R(0)}
= \frac{\omega_0^2}{2\lambda\sqrt{\omega_0^2 - \lambda^2}}
= \frac{\omega_0}{2\lambda}\cdot
\frac{\omega_0}{\sqrt{\omega_0^2 - \lambda^2}}
= Q\,\Bigl(1 - \frac{\lambda^2}{\omega_0^2}\Bigr)^{-1/2}
\geq Q .
\]
Para amortiguamiento débil, el factor de corrección está próximo a $1$:
la \omterm{ex:b1:diffeq:oscillation}{resonancia} multiplica el desplazamiento estático esencialmente
por $Q$.

\textbf{14.} $V(\Omega)^2 = \frac{A^2 u}{g(u)}$ con $u = \Omega^2$. Su
derivada tiene el signo de $g(u) - u\,g'(u) = (\omega_0^2 - u)^2 +
4\lambda^2 u - u\bigl(2(u - \omega_0^2) + 4\lambda^2\bigr) =
(\omega_0^2 - u)^2 + 2u(\omega_0^2 - u) = (\omega_0^2 - u)(\omega_0^2 +
u)$, positivo para $u < \omega_0^2$ y negativo más allá: máximo
estricto exactamente en $\Omega = \omega_0$, sea cual sea el
amortiguamiento. Ahí $\tan\varphi$ estalla con $\varphi \in
\intoo0\pi$: $\varphi = \frac\pi2$, y $x_p'(t) = -R\Omega\sin(\Omega t
- \frac\pi2) = R\Omega\cos(\Omega t)$ está exactamente en fase con la
fuerza: transferencia óptima de potencia.

\textbf{15.} $R = R_{\max}/\sqrt2 \iff g(u) = 2g(u_r) \iff (u -
u_r)^2 = g(u_r) = 4\lambda^2(\omega_0^2 - \lambda^2)$, lo que da
\[
u_\pm = u_r \pm 2\lambda\sqrt{\omega_0^2 - \lambda^2},
\qquad
\Omega_+^2 - \Omega_-^2 = 4\lambda\sqrt{\omega_0^2 - \lambda^2} .
\]
Después, $\Omega_+ - \Omega_- = \frac{\Omega_+^2 -
\Omega_-^2}{\Omega_+ + \Omega_-}$ y, para $\lambda \ll \omega_0$, los
dos $\Omega_\pm \approx \omega_0$: $\Omega_+ - \Omega_- \approx
\frac{4\lambda\omega_0}{2\omega_0} = 2\lambda$, luego
$\frac{\omega_0}{\Omega_+ - \Omega_-} \approx
\frac{\omega_0}{2\lambda} = Q$. Medir el ancho de un pico de
\omterm{ex:b1:diffeq:oscillation}{resonancia} es medir su \omterm{pb:b1:diffeq:1}{factor de calidad}.

\textbf{16.} $Q = 10$; $\Omega_r = \sqrt{1 - 2(0.05)^2} =
\sqrt{0.995} = 0.9975$; $R_{\max} = \frac1{2 \times 0.05
\sqrt{1 - 0.0025}} = \frac1{0.1 \times 0.99875} = 10.01$; respuesta
estática $R(0) = 1$; ancho de banda $\approx 2\lambda = 0.1$. Un pico
alto y estrecho de altura $\approx Q$ sobre una meseta de altura $1$.

\textbf{17.} Si $2\lambda^2 \geq \omega_0^2$, entonces $u_r \leq 0$ y
$g'(u) = 2(u - u_r) > 0$ para todo $u > 0$: $g$ crece estrictamente en
$\intoo0{+\infty}$, luego $R = A/\sqrt g$ decrece estrictamente desde
$R(0) = A/\omega_0^2$: la respuesta es máxima a frecuencia cero y no
hay pico alguno.

\textbf{18.} $\gamma = \iu\Omega$ no es raíz de $r^2 + \omega_0^2$
(pues $\Omega \neq \omega_0$), así que
el~\cref{met:b1:diffeq:particular} con $m = 0$ da
$x_p = \frac{A\cos\Omega t}{\omega_0^2 - \Omega^2}$ (sustitúyase y
compruébese: $-\Omega^2 + \omega_0^2$ multiplicado por el coseno).
Solución general:
\[
x(t) = \frac{A\cos\Omega t}{\omega_0^2 - \Omega^2}
+ \lambda_1\cos\omega_0 t + \mu_1\sin\omega_0 t .
\]

\textbf{19.} $x(0) = 0$ obliga a $\lambda_1 = -\frac A{\omega_0^2 -
\Omega^2}$, y $x'(0) = 0$ obliga a $\mu_1 = 0$:
\[
x(t) = \frac{A}{\omega_0^2 - \Omega^2}\,
\bigl(\cos\Omega t - \cos\omega_0 t\bigr)
= \frac{2A}{\omega_0^2 - \Omega^2}\,
\sin\Bigl(\frac{(\omega_0 - \Omega)t}2\Bigr)
\sin\Bigl(\frac{(\omega_0 + \Omega)t}2\Bigr) ,
\]
por la fórmula de producto $\cos a - \cos b =
2\sin\frac{b - a}2\sin\frac{b + a}2$ aplicada con $a = \Omega t$,
$b = \omega_0 t$.

\textbf{20.} Para $\Omega$ próximo a $\omega_0$, el segundo seno oscila
a la frecuencia rápida $\frac{\omega_0 + \Omega}2 \approx \omega_0$,
mientras que el primero es una envolvente lenta de frecuencia
$\frac{\abs{\omega_0 - \Omega}}2$: la amplitud de la oscilación rápida
crece y decrece con período de envolvente
$\frac{2\pi}{\abs{\omega_0 - \Omega}}$ (dos \omterm{pb:b1:diffeq:1}{pulsaciones} por período de
envolvente), alcanzando máximos
$\frac{2A}{\abs{\omega_0^2 - \Omega^2}}$. Cuando $\Omega \to \omega_0$,
las \omterm{pb:b1:diffeq:1}{pulsaciones} se vuelven a la vez más lentas (período $\to \infty$) y
más altas (amplitud $\to \infty$).

\textbf{21.} Fíjese $t$. Cuando $\Omega \to \omega_0$:
\[
\frac{2A}{\omega_0^2 - \Omega^2}\sin\Bigl(\frac{(\omega_0 -
\Omega)t}2\Bigr)
= \frac{2A}{\omega_0 + \Omega}\cdot
\frac{\sin\bigl(\frac{(\omega_0 - \Omega)t}2\bigr)}
{\omega_0 - \Omega}
\longrightarrow \frac{2A}{2\omega_0}\cdot\frac t2
= \frac{At}{2\omega_0} ,
\]
mientras que $\sin\bigl(\frac{(\omega_0 + \Omega)t}2\bigr) \to
\sin(\omega_0 t)$: el límite es $x_\infty(t) =
\frac{At\sin\omega_0 t}{2\omega_0}$. Comprobación directa: con
$C = \frac A{2\omega_0}$, $x_\infty = Ct\sin\omega_0 t$ cumple
$x_\infty'' = 2C\omega_0\cos\omega_0 t - C\omega_0^2 t
\sin\omega_0 t$, luego $x_\infty'' + \omega_0^2 x_\infty =
2C\omega_0\cos\omega_0 t = A\cos\omega_0 t$, con $x_\infty(0) = 0$ y
$x_\infty'(0) = C\sin 0 + C\omega_0 \cdot 0 \cdot \cos 0 = 0$ --- para
$\omega_0 = A = 1$ esto es exactamente
el~\cref{ex:b1:diffeq:oscillation}. La \omterm{ex:b1:diffeq:oscillation}{resonancia} es la degeneración
de las \omterm{pb:b1:diffeq:1}{pulsaciones}: la primera crecida de la envolvente, estirada hasta
una longitud infinita.

\textbf{22.} Sin amortiguamiento, la amplitud resonante crece
linealmente y sin cota; con amortiguamiento $\lambda > 0$, el
crecimiento se satura en $R_{\max} \approx Q\,\frac A{\omega_0^2}$. El
mecanismo: la disipación retira energía a un ritmo que crece con la
amplitud (pregunta 23), de modo que la acumulación se detiene
exactamente cuando el amortiguamiento quema energía tan deprisa como la
aporta la excitación.

\textbf{23.} Con $x_p' = -R\Omega\sin(\Omega t - \varphi)$, y siendo
los promedios en un período $\langle\cos^2\rangle =
\langle\sin^2\rangle = \frac12$ y $\langle\sin\cos\rangle = 0$:
\[
\langle P_{\mathrm{diss}}\rangle
= 2\lambda\,R^2\Omega^2\,\langle\sin^2(\Omega t -
\varphi)\rangle = \lambda R^2\Omega^2 ;
\]
y, desarrollando $\sin(\Omega t - \varphi) = \sin\Omega t\cos\varphi
- \cos\Omega t\sin\varphi$:
\[
\langle P_{\mathrm{in}}\rangle
= -AR\Omega\,\bigl\langle\cos\Omega t\,\sin(\Omega t -
\varphi)\bigr\rangle
= AR\Omega\,\frac{\sin\varphi}2 .
\]
Como $\sin\varphi = \frac{2\lambda\Omega}{\sqrt D} =
\frac{2\lambda\Omega R}A$, esto vale $\frac{AR\Omega}2 \cdot
\frac{2\lambda\Omega R}A = \lambda R^2\Omega^2$: las potencias
aportada y disipada se equilibran exactamente --- la propiedad que
define un régimen permanente.

\textbf{24.} (i) El teorema de estructura separó cada solución en
régimen permanente más transitorio (preguntas 9--11) y redujo la
unicidad al problema homogéneo. (ii) El método complejo convirtió la
búsqueda de una solución particular en una única división de números
complejos (pregunta 6), leyéndose la amplitud y la fase en un
\omterm{def:b1:complex:field}{módulo} y un argumento. (iii) La curva de \omterm{ex:b1:diffeq:oscillation}{resonancia} es un puro
estudio de función --- una cuadrática en $u = \Omega^2$, su mínimo, sus
\omterm{def:b1:logic:sets}{conjuntos} de nivel --- al estilo del~\cref{ch:b1:functions}
(preguntas 12--17).

\textbf{25.} Libre y amortiguado: seudooscilaciones que decaen, con
vida $\frac1\lambda$ y unas $\frac Q\pi$ oscilaciones. Forzado y
amortiguado: los transitorios mueren y sobrevive un único régimen
permanente sinusoidal a la frecuencia de excitación, con amplitud que
alcanza su pico cerca de $\omega_0$ (altura $\approx Q \times$ la
estática, anchura $\approx \frac{\omega_0}Q$) y fase que barre de $0$ a
$\pi$ pasando por $\frac\pi2$ en $\omega_0$. Libre y sin amortiguar:
oscilación perpetua. Forzado y sin amortiguar: \omterm{pb:b1:diffeq:1}{pulsaciones}, que
degeneran en \omterm{ex:b1:diffeq:oscillation}{resonancia} de crecimiento lineal cuando la sintonía es
exacta. Un solo número adimensional, $Q = \frac{\omega_0}{2\lambda}$,
lo ajusta todo --- altura del pico, ancho de banda y vida del
transitorio son tres lecturas del mismo mando. La continuación:
reescribir $x'' + 2\lambda x' + \omega_0^2x$ como sistema de primer
orden abre los métodos matriciales del~\cref{ch:b1:matrices} y del
volumen del segundo año, y descomponer una excitación periódica
arbitraria en sinusoides (series de Fourier, volumen del tercer año)
convierte el análisis de este problema, a una sola frecuencia, en el
ladrillo universal: resuélvase para cada frecuencia y superpóngase.
\end{solution}
